\subsection{可积性判别法}

定理 5. --- 要使 X 到 Banach 空间 F 中的映射 f 是 p 次幂可积的（ \(1 \leqslant p < +\infty\) ），必须且只需 f 可测且 \(N_{p}(\mathbf{f})\) 有限。

条件是必要的：因为若 \(f \in L_{F}^{p}\)，则存在具有紧支集的连续函数的一个序列 \((\mathbf{g}_{n})\)，它几乎处处收敛于 f（§3, No.~4, 定理 3 的推论 2）；由 No.~4 的定理 2，f 是可测的。

为证明条件是充分的，我们先建立一条引理：

引理 1. --- 设 g 是取值于 F 的一个函数，使得 \(\mathrm{N}_{p}(\mathbf{g}) < +\infty\)（换言之，\(F_{F}^{p}\) 中的一个函数）。点 \(x \in X\)（使得 \(\mathbf{g}(x) \neq 0\) 者）所成之集 A 含于一个可忽略集与一个紧集序列的并。

设 \(A_{n}\) 是点 \(x \in X\)（使得 \(|\mathbf{g}(x)| \geqslant 1/n\) 者）所成之集；A 是诸 \(A_{n}\) 的并，且 \(\varphi_{A_{n}} \leqslant n|\mathbf{g}|\)，于是 \(|\mu|^{*}(A_{n}) \leqslant (nN_{p}(\mathbf{g}))^{p}\)；由此可知 \(A_{n}\) 含于一个可忽略集与一个紧集序列的并（§4, No.~6, 定理 4 的推论 3），因此 A 亦然。

引理证毕后，先考虑 \(\mathbf{f}\) 具有紧支集 \(\mathbf{K}\) 的情形。由 No.~5 的定理 3 的推论 1，存在可测阶梯函数的一个序列 \((\mathbf{g}_n)\)，使得 \(|\mathbf{g}_n(x)| \leqslant |\mathbf{f}(x)|\)（在每一点 \(x \in \mathbf{X}\) 处），且 \(\mathbf{g}_n(x)\) 几乎处处趋于 \(\mathbf{f}(x)\)。而现在，\(\mathbf{g}_n\) 是含于 \(\mathbf{K}\) 的可测集的特征函数的一个线性组合；由于这些集合依据 No.~1 的命题 3 是可积的，故 \(\mathbf{g}_n\) 属于 \(\mathcal{L}_{\mathrm{F}}^p\)。由于 \(\mathrm{N}_p(\mathbf{f}) < +\infty\)，Lebesgue 定理（§3, No.~7, 定理 6）表明 \(\mathbf{f}\) 属于 \(\mathcal{L}_{\mathrm{F}}^p\)。

在一般情形下，由引理 1 可知存在紧集的一个递增序列 \((\mathrm{K}_n)\)，使得 \(\mathbf{f}(x)\) 在诸 \(\mathrm{K}_n\) 的并的余集中几乎处处为零。设 \(\mathbf{f}_n\) 是等于 \(\mathbf{f}(x)\) 于 \(\mathrm{K}_n\) 上、在其他处等于 0 的函数；\(\mathbf{f}_n\) 依据 No.~3 的定理 1 的推论 5 是可测的；由于 \(|\mathbf{f}_n| \leqslant |\mathbf{f}|\)，由论证的第一部分知 \(\mathbf{f}_n\) 属于 \(\mathcal{L}_{\mathrm{F}}^p\)。由于 \(\mathbf{f}(x)\) 几乎处处等于序列 \(\mathbf{f}_n(x)\) 的极限，Lebesgue 定理再次证明 \(\mathbf{f} \in \mathcal{L}_{\mathrm{F}}^p\)，证明至此完成。

应当注意，一个局部可忽略但并非可忽略的函数不是可积的；因此，与一个可积函数局部几乎处处相等的函数未必可积。

推论 1. --- 要使一个集合可积，必须且只需它是可测的且外测度有限。

推论 2. --- 设 \((\mathrm{F}_{n})\) 是拓扑空间的一个序列；对每个指标 n，设 \(f_{n}\) 是 X 到 \(F_{n}\) 中的可测映射，并设 \(f(x) = (f_{n}(x)) \in \mathrm{F} = \prod_{n} \mathrm{F}_{n}\)；最后，设 u 是 \(f(\mathrm{X})\) 到 Banach 空间 G 中的一个连续映射。要使函数 \(\mathbf{g}(x) = \mathbf{u}(f(x))\) 可积，必须且只需 \(\mathrm{N}_{1}(\mathbf{g}) < +\infty\)。

因为 \(\mathbf{g}\) 是可测的（No.~3, 定理 1）。

推论 3. --- 对每个可积函数 f 与每个可测集 A，函数 \(f\varphi_{A}\) 是可积的。

因为由定理 5 与 No.~3 的定理 1 的推论 5 可知 \(\mathbf{f}\varphi_{\mathrm{A}}\) 是可测的，且 \(\mathrm{N}_1(\mathbf{f}\varphi_{\mathrm{A}}) \leqslant \mathrm{N}_1(\mathbf{f})\)。

我们记 \(\int_{\mathrm{A}} \mathbf{f} d\mu = \int \mathbf{f} \varphi_{\mathrm{A}} d\mu\)（或 \(\int_{\mathrm{A}} \mathbf{f} \mu\)），对每个可积函数 \(\mathbf{f}\) 与每个可测集 \(\mathrm{A}\)。我们也用 \(\int_{\mathrm{A}}^{*} f d|\mu|\)（或 \(\int_{\mathrm{A}}^{*} f |\mu|\)）来代替 \(\int^{*} f \varphi_{\mathrm{A}} d|\mu|\)，对每个数值函数（有限或非有限的）\(f \geqslant 0\)（约定 \(f(x) \varphi_{\mathrm{A}}(x) = 0\) 当 \(f(x) = +\infty\) 且 \(\varphi_{\mathrm{A}}(x) = 0\) 时）。

推论 4. --- 对每个序列 \((f_{n})\)（由 X 上 \(\geqslant 0\) 的可测函数组成），

\[
\int^ {*} \left(\sum_ {n} f _ {n}\right) d | \mu | = \sum_ {n} \int^ {*} f _ {n} d | \mu |.\tag{1}
\]

鉴于 §1, No.~3, 定理 3，我们归结为证明：对任意两个可测函数 \(f \geqslant 0\)、\(g \geqslant 0\)（\(X\) 上的），

\[
\int^ {*} (f + g) d | \mu | = \int^ {*} f d | \mu | + \int^ {*} g d | \mu |.\tag{2}
\]

当 \(f\) 与 \(g\) 可积时，这只是积分的可加性。在相反的情形，例如若 \(f\) 不可积，则由定理 5 有 \(\int^{*} f d|\mu| = +\infty\)；更有 \(\int^{*}(f + g) d|\mu| = +\infty\)。

推论 5. --- 对两两不相交的可测集的每个序列 \((\mathbf{A}_{n})\)，

\[
| \mu | ^ {*} \left(\bigcup_ {n} \mathrm{A} _ {n}\right) = \sum_ {n} | \mu | ^ {*} (\mathrm{A} _ {n}).
\]

这由推论 4 应用于诸 \(\varphi_{\mathrm{A}_n}\) 即得。
% semantic-fragments: legacy-input-seam
\relax{}
